% Internal provenance (excluded from the rendered manuscript):
% Implementation and reported results: ff3c7f3f465bdfce9a9a2f0050fbac6db6af82fc.
% Measurements originate from README.md; raw GPU logs were not supplied.
% GPU benchmarks and Lean regeneration/checking were not rerun during drafting.
% Current Q8 limits follow core/ie_core.h; older README bounds are superseded.
% Later model-specific records supersede the README's older coverage statement.
% Source hashes and citation receipts are maintained in supporting JSON files.
% Build: tectonic --keep-logs --keep-intermediates --outdir paper paper/main.tex
\documentclass[10pt]{article}
\usepackage[margin=1in]{geometry}
\usepackage{amsmath,amssymb,booktabs,graphicx}
\usepackage[hidelinks]{hyperref}
\usepackage{microtype}
\title{Bandwidth-Efficient Qwen3.8 Decoding on AMD RDNA4\\with a Verified Integer Core and Multi-Token Prediction}
\author{Author information pending}
\date{Research draft --- October 2026}
\begin{document}
\maketitle
\begin{abstract}
Running a large hybrid language model on one graphics processor requires efficient weight streaming, recurrent-state management, and careful numerical implementation. This paper describes amd-infer, a C inference engine specialized for the Qwen3.8-27B Q8\_0 model and the AMD RDNA4 gfx1201 architecture. The engine combines row-major quantized weight packing, wave-level integer dot products, fused hybrid-attention operations, streaming model loading, and multi token prediction (MTP). A shared core has accompanying Lean proofs for indexing, quantized integer dot products, reduction identities, and activation-plan checking. This verification boundary excludes floating-point attention, compiler correctness, device synchronization, and speculative decoding. Reported measurements show 21.7 tokens/s for ordinary decoding on a Radeon AI PRO R9700, corresponding to approximately 590 GB/s of effective weight bandwidth. With three MTP drafts, recorded throughput reaches 60.1 tokens/s on a coding prompt and 50.6 tokens/s on a text-continuation prompt. The evaluation also reports greedy-token agreement with ordinary decoding and a WikiText-2 perplexity of 5.1804 on four short evaluation chunks. Raw benchmark logs are unavailable, limiting independent assessment of measurement variability. The design illustrates how a narrow verified core can coexist with device-specific optimization while keeping empirical performance and numerical claims separate from proof-backed properties.
\end{abstract}
\noindent\textbf{Keywords:} quantized inference; weight reuse; recurrent state; speculative decoding; memory planning; formal verification

\section{Introduction}
Single-sequence autoregressive decoding repeatedly applies large weight matrices to a small number of activation vectors. For a quantized model that occupies most of a device's memory, repeatedly reading those matrices makes weight traffic a central engineering constraint. Hybrid attention adds another requirement: speculative execution must recover both position-indexed attention caches and recurrent state after rejecting a draft token.

Multi-token prediction (MTP) is the speculation mechanism examined here. The amd-infer engine targets Qwen3.8-27B Q8\_0, a 64-layer hybrid model containing 48 linear-attention layers and 16 full-attention layers. Its quantized model file occupies approximately 28.6 GB. The design combines single-token execution with bounded speculative verification on the AMD RDNA4 architecture.

The system offers three concrete engineering contributions. First, a compact C engine exposes weight layout and wave-level work assignment directly, using eight-bit integer dot products for Q8\_0 matrices. Second, hybrid-attention execution combines a fused recurrent kernel with state snapshots for bounded speculation. Third, a shared integer core connects implementation functions to accompanying Lean proofs, including a checker for activation memory plans. The proof boundary is deliberately explicit: the full decoder is not formally verified.

The evaluation uses previously reported measurements to examine effective weight bandwidth and the throughput gains from multi-token prediction. Numerical checks assess greedy-token agreement and perplexity on a limited text sample.

\section{Background and Related Work}
Quantized language-model inference separates integer representations from floating-point scaling. The Q8\_0 format uses a half-precision scale and 32 signed eight-bit values per block. Activations use their own block scale. Integer products can therefore be accumulated before converting the result to floating point and applying the two scales.

Speculative sampling evaluates draft continuations with a target model and applies an acceptance rule to preserve a target distribution \cite{chen}. The engine examined here implements a narrower greedy variant. Drafts are accepted only when they match the target model's argmax token. No distribution-preserving stochastic sampler is implemented or claimed. The use of the model's multi-token prediction head also differs from deploying an unrelated small draft model.

WikiText was introduced with the pointer-sentinel language-model study \cite{merity}. The recorded quality check uses only four WikiText-2 chunks of 512 tokens each, so it is a short diagnostic rather than a full-corpus or task-level evaluation.

\section{Engine Architecture}
\subsection{Loading and execution graph}
The host parses GGUF (GGML Universal File Format) version 3 metadata and tensors, builds an operator list, and assigns activation buffers to a shared arena. Central processing unit (CPU) and graphics processing unit (GPU) backends consume the same graph and offsets. The GPU host dynamically loads the Heterogeneous-compute Interface for Portability (HIP) runtime, while the kernels are compiled as C for \texttt{amdgcn-amd-amdhsa}, targeting \texttt{gfx1201}. Runtime HIP availability remains necessary even though kernel compilation does not require a conventional HIP C++ source pipeline.

Projection matrices are concatenated when their storage types permit it. In a linear-attention layer, the query, key, value, output-gate, beta, and alpha projections are combined into one matrix operation. Gate and up projections in the feed-forward path are likewise combined. Bias, residual addition, and production of quantized activation copies are fused into selected kernels. These choices reduce intermediate launches and avoid separate quantization passes where supported.

GPU loading is streaming: each repacked matrix is copied to device memory and its temporary host storage is released. Reported loading statistics include 516 uploaded arrays totaling 28.57 GB and approximately 50 seconds of loading time for the 27B model. The Q8\_0 embedding remains quantized and is gathered by row, avoiding expansion of the full embedding matrix to four-byte floating-point values.

\subsection{Activation memory planning}
The planner sorts buffers by size and chooses the lowest available 256-byte-aligned offset that does not conflict with an already placed live buffer. Lifetimes are closed intervals of operator indices. If buffers $i$ and $j$ are simultaneously live, their byte intervals must be disjoint:
\begin{equation}
 [o_i,o_i+s_i)\cap[o_j,o_j+s_j)=\varnothing.
\end{equation}
The candidate plan is checked by \texttt{ie\_plan\_check}; execution aborts if it is rejected. Accompanying proofs establish arena containment, alignment, and non-overlap for accepted plans under the formal model. The greedy construction algorithm itself is not proved optimal or correct. Actual allocation lengths, faithful lifetime annotations, and the caller's use of offsets remain obligations outside the checker theorem.

\section{Quantized Kernels and the Verified Core}
\subsection{Q8 layout and dot products}
A Q8\_0 matrix is repacked into separate row-major arrays of integer words and half-precision scales. Each 32-value block occupies eight 32-bit words in the integer array. For row $r$, block $b$, word $w$, and $n_b$ blocks per row, the destination word index is
\begin{equation}
 i(r,b,w)=8(rn_b+b)+w,\qquad 0\leq w<8.
\end{equation}
The accompanying address and inverse-map laws connect this layout to the source bytes. The current Q8 size predicate requires $r_{\mathrm{count}}<2^{18}$, $n_b<2^{11}$, and $r_{\mathrm{count}}n_b<2^{26}$.

For block $b$, weights are represented by scale $s_b$ and signed integers $q_{bj}$. Activations use $a_{bj}$ and a floating-point scale $d_b$, computed from the maximum absolute activation value divided by 127. The output approximation is
\begin{equation}
 y_r=\sum_{b=0}^{n_b-1}s_{rb}d_b
       \left(\sum_{j=0}^{31}q_{rbj}a_{bj}\right).
\end{equation}
Eight packed dot4 operations implement the inner integer dot product. The core models signed-byte interpretation and accumulation using 32-bit unsigned arithmetic. Integer identities are proved modulo $2^{32}$; they do not establish exact real-valued output equivalence after quantization and floating-point scaling.

\subsection{Wave assignment and memory access}
Each workgroup contains 256 threads, organized as eight 32-lane waves. Wave $v$ in workgroup $g$ owns row $8g+v$. Lane $\ell$ handles blocks $\ell+32t$. The nominal grid has $\lceil R/8\rceil$ workgroups for $R$ rows. Coverage and uniqueness laws establish the intended assignment within the modeled domain.

The kernels stage two blocks of loads before performing their arithmetic. Weight accesses use 128-bit vector loads, and wave-level reductions use data-permutation operations rather than shared-memory permutation. The chosen reduction order is retained across CPU and GPU implementations. This ordering limits one source of numerical discrepancy but cannot make floating-point addition associative.

The GPU launcher pads nominal grids of 253--256 workgroups to 257. A reported microbenchmark takes 17.0 microseconds at 256 groups and 6.5 microseconds at 257 for a 2048-row, 48-block microbenchmark on the R9700. Additional groups exit without producing rows. The cause of this discontinuity is unknown, and the policy should be treated as a device-specific observation rather than a general scheduling rule.

\subsection{Verification boundary}
The C and C++ interfaces share one integer-core header. Its restricted functional C++ view is translated to Lean using cpp-lean and the Edison Design Group (EDG) front end. Proof files cover work assignment, Q4 and Q8 addresses and packing, integer dot products, integer reduction identities, and plan-checker soundness. Q4 proof support exists in the shared core but is not the weight path evaluated in this paper.

The trusted boundary includes the translation from source semantics to Lean, host and GPU compiler correctness, the dot4 instruction matching its modeled behavior, valid caller-supplied lengths and size predicates, and the device memory model. Floating-point normalization, rotary embeddings, softmax, gating, cross-workgroup synchronization, parsing, and speculative decoding are outside these proofs. The analysis relies on the accompanying proof artifacts; independent regeneration and proof checking remain necessary for reproduction.

\section{Hybrid Attention and Multi-Token Prediction}
\subsection{Fused recurrent attention}
The hybrid architecture places three linear-attention layers before each full-attention layer. A linear layer uses a gated delta recurrence with a $128\times128$ state matrix per value head. Its fused kernel applies a four-tap causal convolution and SiLU, normalizes query and key vectors in the Euclidean norm, computes beta and decay gates, updates state, and applies gated root mean square normalization (RMSNorm) to the output.

Let $S_{p-1}$ be the previous key-by-value state, and let $q_p,k_p,v_p$ denote the processed query, key, and value. With $\beta_p=\operatorname{sigmoid}(b_p)$ and $\lambda_p=\exp(\operatorname{softplus}(a_p+\delta_b)A)$, the implemented recurrence is
\begin{align}
 \widetilde S_p&=\lambda_p S_{p-1},\\
 \Delta_p&=\beta_p(v_p-\widetilde S_p^{\mathsf T}k_p),\\
 S_p&=\widetilde S_p+k_p\Delta_p^{\mathsf T},\\
 o_p&=S_p^{\mathsf T}q_p/\sqrt{128}.
\end{align}
A value head selects key head $h\bmod n_k$. The output is normalized, multiplied by a learned weight and $\operatorname{SiLU}(z_p)$, and optionally quantized for the next projection. At position zero, the kernel treats prior state as zero without reading it.

A workgroup owns one value head. State loads are issued early to overlap with convolution computation. A supporting measurement reports approximately 6.9 microseconds per recurrent-kernel invocation for the smaller Qwen3.5-0.8B workload; this is supporting implementation context and is not a measured 27B kernel latency.

\subsection{Full attention}
Full-attention layers allocate position-indexed key--value (KV) caches. Queries and keys receive per-head RMSNorm, and the rotary position transformation operates on the first 64 dimensions with NEOX pairing. Query projections also supply a gate; the attention output is multiplied by its sigmoid.

Attention is split across query heads and 32-position segments. Waves form online-softmax summaries, a workgroup combines local summaries, and the final completing workgroup combines segment maxima, denominators, and weighted values. These operations use shared memory and atomic completion counters. Their synchronization is tested implementation behavior, rather than a theorem of the integer core.

\subsection{Bounded speculation and weight reuse}
The engine supports one to three MTP drafts per verification step. The prediction head first catches up previously accepted positions using the target model's hidden states. It predicts the first draft from a true hidden state and the current token; later drafts use its own predicted hidden state. The target then evaluates the current token and all $D$ drafts in one pass.

For Q8 projections, the multi-token kernel maintains up to four activation accumulators while loading each weight block once. For each activation vector it preserves the single-token arithmetic order. Other operations are executed in token order where necessary; this is limited verification batching, not a general dense matrix-multiplication engine.

If target predictions are $a_0,\ldots,a_D$ and drafts are $d_1,\ldots,d_D$, acceptance chooses the longest prefix satisfying $d_j=a_{j-1}$, then emits the target prediction $a_k$ after $k$ accepted drafts. This yields ordinary greedy tokens provided the batched target computation and restored state match ordinary execution. Reported numerical checks show such token agreement for tested draft depths, but it is not a formal end-to-end equivalence proof.

Recurrent state uses four slots. Within a verification step starting from slot $c$, state after token index $t$ is written to $(c+t)\bmod4$. After accepting $k$ drafts the host retains slot $(c+k)\bmod4$. The convolution input ring has eight slots, preserving the history needed for at most four verified tokens and a four-tap convolution. KV entries are overwritten by position on subsequent execution. This storage protocol makes rollback possible without copying every recurrent matrix after rejection.

\section{Evaluation}
\subsection{Evidence and protocol}
The evaluation uses previously reported measurements on a Radeon AI PRO R9700 with a nominal memory bandwidth of 640 GB/s. The implementation targets gfx1201 and uses a ROCm 7.1/7.14 runtime environment. Results cover ordinary decoding, MTP throughput, and a short WikiText-2 perplexity check. Raw per-run timing logs, exact prompt strings, model-file hashes, and a complete machine configuration are unavailable; these omissions limit statistical analysis and independent reproduction.

\subsection{Ordinary decoding and bandwidth}
Table~\ref{tab:ordinary} reports ordinary Qwen3.8-27B decoding. Each length was reportedly run three times with less than 0.1\% variation. Individual samples and the definition of this variation are unavailable, so no confidence interval is estimated. Effective bandwidth counts approximately 27.2 GB of weights per generated token. It is a workload-derived ratio, not a hardware-counter measurement of all memory traffic.
\begin{table}[ht]
\centering
\caption{Recorded ordinary Qwen3.8-27B Q8\_0 decoding on R9700.}
\label{tab:ordinary}
\begin{tabular}{lrr}
\toprule
Generated tokens & Milliseconds/token & Tokens/s\\
\midrule
64 & 46.01 & 21.7\\
512 & 46.17 & 21.7\\
\bottomrule
\end{tabular}
\end{table}
Using the 64-token latency,
\begin{equation}
 B_{\mathrm{eff}}=\frac{27.2\ \mathrm{GB}}{0.04601\ \mathrm{s}}
 \approx591\ \mathrm{GB/s},\qquad
 \frac{B_{\mathrm{eff}}}{640\ \mathrm{GB/s}}\approx92.4\%.
\end{equation}
Rounded to the precision of the reported measurements, effective bandwidth is approximately 590 GB/s, or 92\% of nominal bandwidth. Separate Q8 shape microbenchmarks for the 27B projections report 589--634 GB/s. Their proximity to the stated bandwidth ceiling is consistent with weight traffic dominating these operations, but does not establish a counter-based bottleneck diagnosis for every kernel.

\subsection{MTP throughput}
Table~\ref{tab:mtp} reports 300-token generation runs. The prompts cover coding, text continuation, and chat; exact prompt strings are unavailable. Acceptance is the number of accepted drafts divided by the number of attempted drafts. Speedups in the table are calculated against the recorded 21.7 tokens/s ordinary rate.
\begin{table}[ht]
\centering
\caption{Recorded Qwen3.8-27B MTP results; speedups are arithmetic ratios.}
\label{tab:mtp}
\begin{tabular}{llrrr}
\toprule
Prompt category & Drafts & Tokens/s & Acceptance & Speedup\\
\midrule
Coding & 0 & 21.7 & --- & 1.00\\
Coding & 1 & 38.4 & 98.7\% & 1.77\\
Coding & 2 & 51.3 & 97.1\% & 2.36\\
Coding & 3 & 60.1 & 93.2\% & 2.77\\
Text continuation & 0 & 21.7 & --- & 1.00\\
Text continuation & 1 & 36.0 & 86.9\% & 1.66\\
Text continuation & 2 & 45.6 & 81.6\% & 2.10\\
Text continuation & 3 & 50.6 & 74.2\% & 2.33\\
Chat & 3 & 47.5 & 67.3\% & 2.19\\
\bottomrule
\end{tabular}
\end{table}
The chat ratio uses the same ordinary rate; a separate matched chat baseline is not reported. Acceptance falls as draft depth increases, yet measured throughput rises through three drafts in the two reported depth sweeps. This behavior is consistent with amortizing the large target weight read over more useful tokens. The data do not justify extrapolation to other prompts or larger draft depths.

As a simple accounting model, if $K$ drafts are accepted, a step produces approximately $1+K$ useful tokens and costs target verification plus draft work. Thus
\begin{equation}
 \mathrm{throughput}\approx\frac{1+\mathbb E[K]}{t_{\mathrm{verify}}+t_{\mathrm{draft}}}.
\end{equation}
A reported timing estimate is roughly 63 ms per step at three drafts, compared with about 46 ms per ordinary token, and attributes additional work mainly to the MTP and vocabulary-output projections. This attribution is a hypothesis rather than a measured decomposition. Each draft reads approximately 0.45 GB of prediction-head weights and 1.35 GB of output-layer weights.

\subsection{Numerical checks}
The recorded Qwen3.8-27B perplexity is 5.1804 on four WikiText-2 chunks of 512 tokens. Only the latter 255 predicted tokens of each chunk are scored. If their negative log-probabilities are $\ell_i$, perplexity is $\exp(\sum_i\ell_i/N)$, giving $N=1020$ scored predictions across the four chunks. This small sample does not establish full-corpus perplexity or downstream task quality.

Reported numerical checks show identical greedy tokens between speculative and ordinary execution at draft depths one through three for tested 27B cases. CPU dot-product tests and small synthetic-model comparisons exercise shared integer and floating-point paths, but do not independently validate the entire 27B recurrent computation. Per-position 27B logits and exact-token transcripts are unavailable for independent analysis.

\section{Limitations and Reproducibility}
The available evidence supports an implementation account and a bounded report of prior measurements. It cannot support statistical significance, broad prompt-level generalization, or superiority across inference frameworks. No cross-framework comparison is included in this paper. Backend versions, power settings, clocks, warm-up behavior, exact prompts, and run-level timing samples must be recorded before making stronger empirical claims.

The evaluated configuration uses approximately 451 MB of MTP tensors extracted from a separate Q8\_0 model artifact. Compatibility between the target model and prediction head is consequential for reproduction. Exact artifact identifiers, hashes, and compatibility metadata are unavailable, preventing independent confirmation of this pairing.

Proof-backed properties concern a restricted integer core. The formal tools, compiler chain, memory allocation, runtime state transitions, and floating-point arithmetic remain trusted or empirically checked. Particularly consequential unproved mechanisms include atomic completion in segmented attention and restoring recurrent state after speculative rejection. The reported checks do not establish complete model or context coverage.

A reproduction should first regenerate and check Lean proofs with the documented translator and tool versions, then run integer and synthetic-model checks, record model hashes, and capture device/runtime metadata. Ordinary and speculative runs should use identical prompt tokens and stopping rules, retain output-token sequences, and log individual timings and accepted-prefix lengths. GPU profiling should be performed separately from throughput measurements because per-kernel instrumentation changes execution overhead.

\section{Conclusion}
amd-infer combines quantized wave-level computation, fused hybrid attention, streaming loading, checked activation planning, and bounded MTP verification for Qwen3.8-27B decoding. Recorded ordinary decoding reaches approximately 92\% of the nominal memory-bandwidth ceiling. Three-draft MTP reaches 60.1 tokens/s on the recorded coding prompt and 50.6 tokens/s on text continuation. These observations suggest that weight reuse can increase useful decoding throughput even when ordinary execution already streams weights efficiently. The accompanying integer proofs provide a narrow foundation for layout, dot products, work assignment, and plan checking; wider correctness and performance claims require additional empirical and formal evidence.

\section*{Declarations}
\paragraph{Data and code availability.} Source code and accompanying Lean proof artifacts are available at \url{https://github.com/hanw/amd.c}. The evaluation uses previously reported measurements. Raw GPU timing logs and model artifacts are not distributed with the manuscript.
\paragraph{Ethics declaration.} This manuscript describes software and recorded machine benchmarks. No human-participant study was performed in preparing it. Model and dataset licensing require confirmation by the authors before redistribution.
\paragraph{Author contributions.} Contributor identities and Contributor Roles Taxonomy (CRediT) assignments require author confirmation. This draft does not attribute implementation or proof work to an unconfirmed person.
\paragraph{Conflict of interest.} Author declarations have not been supplied and remain to be confirmed.
\paragraph{Funding.} Funding information has not been supplied and remains to be confirmed.
\paragraph{AI assistance.} Codex assisted with manuscript drafting and editing. The authors retain responsibility for the technical claims, source attribution, and final manuscript.

% Internal artifact map: README.md (measurements); core/ie_core.h (integer core);
% laws/ie_laws.cpp and proofs/ (formal artifacts); src/model.c and src/plan.c
% (graph and arena); src/hip.c (dispatch); src/main.c (acceptance and rollback);
% kernels/ie_kernels.c (device kernels); tools/gemv_bench.c (microbenchmarks).
% Reproduction commands: make test; make proofs (cpp-lean, EDG, Lean 4.34.0).

\begin{thebibliography}{9}
\bibitem{chen} C. Chen, S. Borgeaud, G. Irving, J.-B. Lespiau, L. Sifre, and J. Jumper, ``Accelerating Large Language Model Decoding with Speculative Sampling,'' arXiv:2302.01318, 2023. \url{https://doi.org/10.48550/arXiv.2302.01318}
\bibitem{merity} S. Merity, C. Xiong, J. Bradbury, and R. Socher, ``Pointer Sentinel Mixture Models,'' arXiv:1609.07843, 2016. \url{https://doi.org/10.48550/arXiv.1609.07843}
\end{thebibliography}
\end{document}
